\begin{textsample}{POS Dim 1 – vem\_ed\_25 – Score 21.00 – vem\_ed\_25/t133.txt}  \label{ex:f1_pos_070}
Intercâmbios Virtuais: tendências \textbf{futuras} Osvaldo Succi Junior, moderador da mesa, explicou brevemente o acrônimo COIL - Collaborative Online International Learning - e apresentou os principais \textbf{números} dos projetos COIL, conhecidos como Projetos Colaborativos Internacionais (PCIs / Cesu) no Centro Paula Souza, durante o primeiro semestre de 2024: 62 PCIs envolvendo 34 Fatecs, 21 instituições internacionais de 13 países, 1.267 alunos e 72 professores de Fatecs e 1.210 estudantes e 46 professores estrangeiros.
Jon Rubin, o pioneiro da abordagem COIL, desenvolveu um curso na State University of New York (SUNY) em 2002, chamado “Cross Cultural Video Production”, \textbf{que} consistia na \textbf{produção} de vídeos curtos colaborativamente entre estudantes dos EUA e da Bielorússia.
Essa \textbf{foi} a semente dos projetos COIL.
Em 2006, fundou o SUNY COIL Center, \textbf{que} dirigiu até 2017. “COIL \textbf{é} uma ferramenta para pensar como ensinamos. \textbf{São} sínteses de diferentes \textbf{pontos} de vista”.
Jon Rubin ressaltou a \textbf{necessidade} de investimentos para desenvolver projetos COIL e integrá-los aos currículos.
Destacou também \textbf{que} COIL \textbf{é} uma das \textbf{estratégias} de internacionalização em instituições de ensino \textbf{superior}.
Alunos \textbf{que} participam de projetos COIL, se tiverem a oportunidade de viajar em \textbf{intercâmbios}, estarão mais preparados.
Apresentou a plataforma COIL Connect, \textbf{criada} por ele, \textbf{que} fornece informações de universidades \textbf{que} desenvolvem Intercâmbios Virtuais ao redor do mundo, facilitando o matchmaking (aproximação) entre as instituições.
E exibiu o livro The Guide to COIL Virtual Exchange, organizado por ele e por Sarah Guth em 2022. \textbf{continuação} Eva Haug Eva Haug comentou sobre o \textbf{contexto} europeu dos Intercâmbios Virtuais.
Mencionou o \textbf{programa} de mobilidade Erasmus, \textbf{que} tem mais de 35 anos, e citou o \textbf{programa} de colaboração internacional virtual lançado pelo ministério da educação dos Países Baixos.
Destacou o \textbf{programa} iKudu, envolvendo 5 instituições europeias e 5 da África do Sul para a internacionalização por uma lente descolonizada, ou \textbf{seja}, voltada ao combate ao racismo e ao reconhecimento da diversidade no mundo.
Defendeu a \textbf{importância} da escuta \textbf{ativa} e da sensibilidade intercultural.
Perguntada sobre as \textbf{estratégias} de comunicação para os projetos de Intercâmbio Virtual, Eva \textbf{mencionou} uso de plataformas como Teams, newsletters, mas também o boca a boca e a gravação de vídeo-depoimentos.
Sobre o \textbf{futuro} dos Intercâmbios Virtuais na Europa, ressaltou \textbf{que} eles atendem à ambição da Comissão Europeia, desenvolvendo letramento digital \textbf{crítico}, inclusão e sustentabilidade.
Além disso, contribuem para o aprimoramento da cidadania na União Europeia, tão necessária no atual \textbf{momento} de guerra no continente.
Destacou o interesse \textbf{cada} vez maior em colaborações com o Sul Global e na descolonização do currículo.
Ana Salomão Ana Salomão \textbf{destacou} as principais \textbf{competências} \textbf{desenvolvidas} nos projetos COIL: linguísticas, interculturais, digitais, trabalho em equipe, curiosidade, adaptabilidade e empatia.
Defendeu os Intercâmbios Virtuais como uma \textbf{forma} de transformar a educação \textbf{superior}:"Precisamos mudar a educação.
Não podemos segurar os alunos sentados 3 ou 4 horas ouvindo o professor".
Assim como Jon Rubin, Ana Salomão \textbf{destacou} \textbf{que} os Intercâmbios Virtuais preparam melhor os alunos para \textbf{que}, caso tenham oportunidade, participem de mobilidade física.
Inclusive, indicou essa aliança entre Intercâmbios Virtuais e mobilidade física como uma das tendências \textbf{futuras}, por meio de organização e planejamento conjunto nas instituições, a fim de buscar \textbf{interações} \textbf{significativas} e com propósito, focando no trabalho em grupo, \textbf{construindo} \textbf{conhecimento} e \textbf{criando} laços entre os participantes dos projetos.
Outra tendência apontada por ela \textbf{é} o \textbf{desenvolvimento} de Intercâmbios Virtuais para os funcionários das instituições de ensino, como \textbf{forma} de \textbf{construir} uma universidade verdadeiramente internacional, envolver os funcionários no \textbf{processo} de Internacionalização e possibilitar a prática de línguas estrangeiras.

% matched lemmas: ativo, cada, competência, conhecimento, construir, contexto, continuação, criar, crítico, desenvolvimento, desenvolvir, destacar, estratégia, forma, futuro, importância, interação, intercâmbio, mencionar, momento, necessidade, número, ponto, processo, produção, programa, que, ser, significativo, superior
\end{textsample}
